\documentclass[11pt,a4paper]{article}

\usepackage[utf8]{inputenc}
\usepackage[T1]{fontenc}
\usepackage{lmodern}
\usepackage{amsmath,amssymb,amsfonts,mathtools}
\usepackage[margin=2.5cm]{geometry}
\usepackage{hyperref}
\usepackage{booktabs}
\usepackage{xcolor}
\usepackage{fancyhdr}
\usepackage{enumitem}
\usepackage{array}
\usepackage{longtable}

\newcommand{\dd}{\mathrm{d}}

\definecolor{statusamber}{RGB}{155,96,0}
\definecolor{sctblue}{RGB}{0,51,153}

\pagestyle{fancy}
\fancyhf{}
\fancyhead[L]{\small\textsc{SCT Theory --- Prediction Card}}
\fancyhead[R]{\small\textsc{SP-01}}
\fancyfoot[C]{\thepage}

\hypersetup{
  colorlinks=true,
  linkcolor=blue!60!black,
  citecolor=green!40!black,
  urlcolor=blue!60!black,
  pdftitle={SCT Prediction Probe: Scalar Mode Decoupling},
  pdfauthor={David Alfyorov},
}

\title{%
  \textbf{SCT Prediction Probe 01:}\\[6pt]
  \Large Scalar-Mode Decoupling at Conformal Coupling\\[4pt]
  \large A Prediction-Style Document for the Curvature-Squared Sector
}
\author{David Alfyorov}
\date{March 2026}

\begin{document}
\maketitle


\begin{center}
\small\textit{Credits: Aliaksandr Samatyia contributed research-assistance and workflow support.}
\end{center}

\begin{center}
  \colorbox{statusamber!15}{\parbox{0.68\textwidth}{\centering
    \textbf{\color{statusamber} STATUS: CONDITIONAL / INDIRECTLY TESTABLE}\\[2pt]
    \small The structural coefficient relation is fixed by the verified
    spectral sum, but direct measurement of the scalar channel remains
    difficult with present data.
  }}
\end{center}

\vspace{1em}

\section{Source}
\label{sec:source}

\begin{itemize}[leftmargin=2em]
  \item \textbf{Derived from:} the verified curvature-squared sector of the
  spectral action after summing the Standard Model field content with the
  correct multiplicities.
  \item \textbf{Relevant local coefficients:}
  \begin{align}
    \alpha_C &= \frac{13}{120},
    \label{eq:alphaC}\\[4pt]
    \alpha_R(\xi) &= 2\left(\xi - \frac{1}{6}\right)^2,
    \label{eq:alphaR}\\[4pt]
    3c_1 + c_2 &= 6\left(\xi - \frac{1}{6}\right)^2.
    \label{eq:3c1c2}
  \end{align}
  \item \textbf{Interpretive step:} in local quadratic gravity, the
  combination $3c_1 + c_2$ controls the scalar spin-0 response in the
  linearized channel.  Therefore \eqref{eq:3c1c2} is the relevant quantity
  for deciding whether the curvature-squared sector introduces an additional
  scalar gravitational mode.
  \item \textbf{Claim type:} this is not a free EFT parametrization.  It is a
  restricted spectral relation among couplings that generic EFT would leave
  unrelated.
\end{itemize}

\section{Observable Quantity}
\label{sec:observable}

\subsection{The prediction}

The SCT prediction is that the scalar-channel coefficient is not arbitrary:
\begin{equation}
  \boxed{
    3c_1 + c_2
    = 6\left(\xi - \frac{1}{6}\right)^2.
  }
  \label{eq:main-prediction}
\end{equation}
Two special cases are immediate:
\begin{itemize}
  \item \textbf{Conformal coupling} $(\xi = 1/6)$:
  \begin{equation}
    3c_1 + c_2 = 0.
    \label{eq:conformal-decoupling}
  \end{equation}
  At the level of the local curvature-squared sector, the extra scalar mode
  decouples.
  \item \textbf{Non-conformal coupling} $(\xi \neq 1/6)$:
  the scalar response returns quadratically with the distance from
  conformality.
\end{itemize}

\subsection{The effective action}

The relevant local effective action is
\begin{equation}
  S_{\mathrm{grav}}
  = \int \dd^4x\,\sqrt{g}\,
  \left[
    \frac{1}{16\pi G}R
    + c_1R^2
    + c_2R_{\mu\nu}R^{\mu\nu}
    + \cdots
  \right].
  \label{eq:effective-action}
\end{equation}
In generic EFT, $c_1$ and $c_2$ are independent Wilson coefficients.
In SCT, once the verified spectral input is inserted, their scalar-channel
combination is forced onto the one-parameter curve \eqref{eq:main-prediction}.

\subsection{Immediate reading}

This prediction should be read as a restriction of parameter space rather
than as a stand-alone low-energy number.  It says:
\begin{quote}
\emph{If the scalar gravitational response sourced by the local
curvature-squared sector is measured, it cannot take an arbitrary value in
SCT.  It must lie on a parabola centered at $\xi = 1/6$.}
\end{quote}

\section{SCT vs.\ Standard EFT and Other Approaches}
\label{sec:comparison}

\begin{center}
\renewcommand{\arraystretch}{1.25}
\begin{tabular}{>{\raggedright\arraybackslash}p{3.1cm} >{\raggedright\arraybackslash}p{4.1cm} >{\raggedright\arraybackslash}p{5.3cm}}
\toprule
\textbf{Approach} & \textbf{Input data} & \textbf{Consequence for scalar mode} \\
\midrule
Generic EFT of gravity & $c_1$, $c_2$ independent & Scalar channel unconstrained until measured \\
Stelle-type quadratic gravity & Bare couplings chosen freely & Extra scalar mode present unless tuned away \\
SCT, general $\xi$ & $3c_1+c_2 = 6(\xi-1/6)^2$ & Scalar channel fixed to a parabola around conformality \\
SCT, $\xi = 1/6$ & $3c_1+c_2 = 0$ & Local scalar mode decouples without extra tuning \\
\bottomrule
\end{tabular}
\end{center}

The distinction is structural.  In a generic quadratic theory, scalar-mode
decoupling is an accidental tuning relation among independent couplings.  In
SCT it is the natural image of the spectral sum at the conformal point.

\section{Physical Observables Sensitive to the Scalar Channel}
\label{sec:observables}

\subsection{Short-distance corrections to the Newtonian potential}

If a scalar mode is present, the weak-field static potential generically
acquires an extra Yukawa-like contribution.  If the scalar mode decouples,
the tensor channel dominates and the short-distance correction pattern is
different.  Precision short-range gravity experiments therefore constrain the
existence of the scalar response, though the exact constraint depends on the
scale suppressing the curvature-squared operators.

\subsection{Linearized gravitational-wave propagation}

A scalar gravitational degree of freedom may manifest as an additional
polarization or modified scalar-channel propagation in compact-binary
signals.  If SCT sits exactly at $\xi = 1/6$, the local curvature-squared
sector predicts no such scalar contribution at leading order in that channel.

\subsection{Cosmological scalar perturbations}

An additional scalar gravitational response affects the relation between the
metric potentials in linear cosmological perturbation theory and may alter
growth of structure.  A decoupled scalar channel instead places the theory in
a more restrictive subset of quadratic-gravity phenomenology.

\subsection{Black-hole ringdown and quasi-normal spectra}

The presence or absence of an extra scalar sector can, in principle, modify
the quasi-normal spectrum of compact objects.  This is not yet a clean direct
test of the local coefficient relation by itself, but it is the kind of
observable where a scalar/tensor decomposition may eventually become
accessible.

\section{Energy Scale and Regime}
\label{sec:scale}

This prediction belongs to the weak-field, curvature-squared regime.  It
should not be read as a statement about fully nonlinear strong-field dynamics
without additional analysis.  The proper hierarchy is:
\begin{enumerate}[label=\arabic*.]
  \item the local spectral sum fixes the coefficient relation,
  \item the linearized field equations translate that relation into a
  propagator statement,
  \item observables are then computed from the resulting weak-field response.
\end{enumerate}

Thus the prediction is sharp at the theory level, but experimentally useful
only once the relevant kinematic regime and scale are specified.

\section{Required Experimental Precision}
\label{sec:precision}

\subsection{Direct determination}

A direct extraction of $(c_1,c_2)$ separately would be ideal, but is not
currently realistic.  Existing data do not isolate these Wilson coefficients
cleanly in a model-independent manner.

\subsection{Indirect constraints}

A more realistic route is to constrain effective combinations that control
the scalar sector of the weak-field response.  In that sense
\eqref{eq:main-prediction} is already useful: any eventual measurement of the
scalar-channel coefficient must either land on the SCT parabola or falsify
it.

\subsection{Current empirical status}

At present, the best statement is negative rather than positive: there is no
clean empirical determination of the scalar-channel coefficient, but there is
also no evidence that forces it away from the conformal locus.  The
prediction is therefore best classified as open and indirectly testable.

\section{Special Property: Symmetry Around Conformal Coupling}
\label{sec:special}

The scalar-channel coefficient depends only on
\begin{equation}
  \left(\xi - \frac{1}{6}\right)^2.
  \label{eq:symmetric-square}
\end{equation}
Therefore it is invariant under
\begin{equation}
  \xi - \frac{1}{6}
  \longmapsto
  -\left(\xi - \frac{1}{6}\right).
  \label{eq:reflection}
\end{equation}
This means that the local scalar response fixes the distance from the
conformal point, but not the sign of that deviation.  Geometrically, the
conformal coupling is the unique minimum of the scalar channel.

That symmetry is not a decorative observation.  It tells us that the theory
organizes the scalar sector around conformality rather than around minimal
coupling.

\section{Status}
\label{sec:status}

\begin{longtable}{p{4.4cm}p{9.1cm}}
\toprule
\textbf{Item} & \textbf{Status} \\
\midrule
Local coefficient relation & Fixed by the verified spectral sum \\
Conformal decoupling point & Present already in the local coefficient identity \\
Nonlocal completion & Requires the corresponding linearized analysis \\
Direct measurement & Not available at present \\
Indirect constraints & Conceptually available through weak-field observables \\
Overall classification & Genuine structural prediction, not yet directly tested \\
\bottomrule
\end{longtable}

\section{What Would Falsify It}
\label{sec:falsify}

The prediction fails if any of the following occurs:
\begin{enumerate}[label=\arabic*.]
  \item the underlying spectral derivation of the local coefficients is shown
  to be wrong;
  \item a future extraction of the scalar-channel coefficient is incompatible
  with \eqref{eq:main-prediction};
  \item the physically relevant weak-field completion necessarily reintroduces
  an unsuppressed scalar channel even at $\xi = 1/6$.
\end{enumerate}

This last point is especially important.  The present document makes a local
prediction.  A stronger global propagator statement belongs to the linearized
field-equation analysis, not to the prediction card by itself.

\section{Relation to Other SCT Predictions}
\label{sec:relation}

This prediction occupies an intermediate place in the SCT program.

\begin{itemize}
  \item It is more concrete than a foundational postulate because it constrains
  a specific physical sector.
  \item It is less final than a fully phenomenological result because
  translating the scalar-channel coefficient into present-day data still
  requires the full weak-field and nonlocal analysis.
  \item It links directly to the broader SCT claim that spectral geometry
  restricts the gravitational parameter space more strongly than generic EFT.
\end{itemize}

\section{Summary}
\label{sec:summary}

The local curvature-squared sector of SCT predicts
\[
  3c_1 + c_2 = 6\left(\xi - \frac{1}{6}\right)^2.
\]
Hence the conformal point $\xi = 1/6$ is the point of local scalar-mode
decoupling, and any departure from that point restores the scalar channel
quadratically.

This is a proper prediction in the strict sense: the theory removes generic
freedom rather than adding new adjustable parameters.  It does not yet amount
to a direct empirical detection, but it does mark a clean falsifiable target
for any future extraction of the scalar gravitational sector.

\begin{thebibliography}{99}

\bibitem{ChamseddineConnes1997}
A.~H.~Chamseddine and A.~Connes,
``The spectral action principle,''
\emph{Commun. Math. Phys.} \textbf{186} (1997) 731--750,
arXiv:\href{https://arxiv.org/abs/hep-th/9606001}{hep-th/9606001}.

\bibitem{ChamseddineConnesMarcolli2007}
A.~H.~Chamseddine, A.~Connes and M.~Marcolli,
``Gravity and the standard model with neutrino mixing,''
\emph{Adv. Theor. Math. Phys.} \textbf{11} (2007) 991--1089,
arXiv:\href{https://arxiv.org/abs/hep-th/0610241}{hep-th/0610241}.

\bibitem{Donoghue1994}
J.~F.~Donoghue,
``General relativity as an effective field theory: The leading quantum
corrections,''
\emph{Phys. Rev. D} \textbf{50} (1994) 3874--3888,
arXiv:\href{https://arxiv.org/abs/gr-qc/9405057}{gr-qc/9405057}.

\bibitem{Stelle1977}
K.~S.~Stelle,
``Renormalization of higher-derivative quantum gravity,''
\emph{Phys. Rev. D} \textbf{16} (1977) 953--969.

\bibitem{Vassilevich2003Pred}
D.~V.~Vassilevich,
``Heat kernel expansion: User's manual,''
\emph{Phys. Rept.} \textbf{388} (2003) 279--360,
arXiv:\href{https://arxiv.org/abs/hep-th/0306138}{hep-th/0306138}.

\end{thebibliography}

\end{document}
